\section*{अध्याय \ref{ch:b3:heterocycles} --- विषमचक्रीय रसायन}

\begin{solution}{exo:b3:heterocycles:1}
ऑक्साज़ोल: O दो देता है (एक एकाकी युग्म $\pi$ निकाय में, एक तल में), N एक,
तीनों कार्बन तीन: 6। थायाज़ोल: S के साथ वही: 6। पिरिमिडीन: छहों परमाणुओं में से
प्रत्येक से एक, दोनों नाइट्रोजन एकाकी युग्म तल में: 6। प्यूरीन: नौ परमाणु,
N--H नाइट्रोजन दो देता है और शेष एक-एक: 10, एक ह्युकेल संख्या
($4n + 2$, $n = 2$)।
\end{solution}

\begin{solution}{exo:b3:heterocycles:2}
पिपेरिडीन (11.1) $>$ DMAP (9.6) $>$ इमिडाज़ोल (7.0) $>$ पिरिडीन (5.2) $\gg$ पिरोल (लगभग $-4$,
कार्बन पर प्रोटॉनित)।
\end{solution}

\begin{solution}{exo:b3:heterocycles:3}
2-ब्रोमोपिरोल (पिरोल इतना अभिक्रियाशील है कि उसका सहज ही बहुब्रोमीनन हो जाता है);
2-ब्रोमोफ़्यूरान; 3-ब्रोमोइंडोल।
\end{solution}

\begin{solution}{exo:b3:heterocycles:4}
2,5-डाइमेथिलफ़्यूरान; 1,2,5-ट्राइमेथिलपिरोल; 2,5-डाइमेथिलथायोफ़ीन।
\end{solution}

\begin{solution}{exo:b3:heterocycles:5}
2-मेथॉक्सीपिरिडीन। C2 पर मेथॉक्साइड का योग \omterm{def:b3:heterocycles:snar}{माइज़नहाइमर संकुल} के ऋण आवेश को
नाइट्रोजन पर रखता है; C3 पर वह केवल कार्बनों पर फैल सकता है, अतः
मध्यवर्ती, और उससे पहले की संक्रमण अवस्था, ऊर्जा में कहीं ऊँचे होते हैं।
\end{solution}

\begin{solution}{exo:b3:heterocycles:6}
2-ऐमीनोपिरिडीन (उसके सोडियम लवण के रूप में, फिर उपचार-पश्चात् प्रक्रिया में प्रोटॉनित)। ऐमाइड आयन
C2 पर जुड़ता है, जिससे एक ऋणायनी $\sigma$-योगज बनता है जिसमें C2 पर H और $\ce{NH2}$ हैं और ऋण
आवेश वलय-नाइट्रोजन पर है; यह हाइड्राइड खोता है, जो
ऐमीनो समूह से एक प्रोटॉन लेकर $\ce{H2}$ के रूप में निकल जाता है।
\end{solution}

\begin{solution}{exo:b3:heterocycles:7}
N-ऑक्साइड का ऑक्सीजन एक एकाकी युग्म वापस वलय में देता है: अनुनादी
संरचनाएँ C2, C4 और C6 पर ऋण आवेश रखती हैं। इलेक्ट्रॉनरागी C4 पर आक्रमण करते हैं (C2
धनावेशित नाइट्रोजन के निकट है)। फिर ऑक्सीजन को
फ़ॉस्फ़ोरस ट्राइक्लोराइड से हटाया जाता है, जिससे 4-नाइट्रोपिरिडीन मिलता है।
\end{solution}

\begin{solution}{exo:b3:heterocycles:8}
2-हाइड्रॉक्सीपिरिडीन (C2 पर एक OH) और 2-पिरिडोन (N--H और C=O)। पिरिडोन वलय में छह
$\pi$ इलेक्ट्रॉन बनाए रखता है: उसकी ज़्विटर-आयनी अनुनादी संरचना, $\mathrm{O^-}$ और अपना एकाकी युग्म साझा करने वाले
$\mathrm{N^+}$ के साथ, एक पिरिडीनियम-ओलेट है; यह एक ऐमाइड है जिसका
नाइट्रोजन एकाकी युग्म षट्क को पूरा करता है, और प्रबल C=O तथा N--H
हाइड्रोजन आबंध जल में इसका पक्ष लेते हैं।
\end{solution}

\begin{solution}{exo:b3:heterocycles:9}
2-नाइट्रोबेंज़ैल्डिहाइड, मेथिल ऐसीटोऐसीटेट के दो तुल्यांक, और अमोनिया:
कैल्सियम-चैनल अवरोधक निफ़ेडिपीन।
\end{solution}

\begin{solution}{exo:b3:heterocycles:10}
$\ce{CH2}$ समूह की ओर वाला ईन-हाइड्रैज़ीन (आंतरिक, अधिक प्रतिस्थापित)
2,3-डाइमेथिलइंडोल देता है; मेथिल समूह की ओर वाला (अंतस्थ) 2-एथिलइंडोल देता है।
दुर्बल अम्ल मुख्यतः पहला देते हैं, अधिक स्थायी ईन-हाइड्रैज़ीन के माध्यम से;
प्रबलतर अम्ल 2-एथिलइंडोल का अनुपात बढ़ाते हैं।
\end{solution}

\begin{solution}{exo:b3:heterocycles:11}
$\mathrm S_\mathrm NAr$ में दर-निर्धारक चरण योग है, जो C--X
आबंध को नहीं तोड़ता; अधिक विद्युतऋणी फ़्लुओरीन ऋणायनी संक्रमण अवस्था
और \omterm{def:b3:heterocycles:snar}{माइज़नहाइमर संकुल} को अधिक स्थायी करता है। $\mathrm S_\mathrm N2$ में C--X आबंध
संक्रमण अवस्था में टूटता है, और दुर्बल C--I आबंध सबसे तेज़ी से निकलता है।
\end{solution}

\begin{solution}{exo:b3:heterocycles:12}
पिरिडीन: 2\,:\,1\,:\,2 अनुपात में तीन संकेत, पहला \qty{8.6}{ppm} के निकट। पिरोल: 6.7 और \qty{6.2}{ppm} पर
दो समान संकेत और \qty{8.1}{ppm} के निकट एक चौड़ा N--H। फ़्यूरान: 7.4 और \qty{6.4}{ppm} पर दो समान संकेत,
कोई N--H नहीं। थायोफ़ीन: 7.33 और \qty{7.12}{ppm} पर दो समान संकेत, एक-दूसरे के
निकट, फ़्यूरान के संकेतों के बीच।
\end{solution}

\begin{solution}{pb:b3:heterocycles:1}
\textbf{1.} $\ce{C6H5NHNH2 + (CH3)2CO -> C6H5NHN=C(CH3)2 + H2O}$।

\textbf{2.} ऐमीन नाइट्रोजन कार्बोनिल कार्बन पर जुड़ता है, फिर जल निकलता है, जैसे
इमीन बनने में; अम्ल कार्बोनिल समूह को सक्रिय करता है और योगज के OH को
निर्गामी समूह में बदल देता है।

\textbf{3.} अंतस्थ $\ce{NH2}$: दूसरे नाइट्रोजन का एकाकी युग्म फ़ेनिल वलय के साथ संयुग्मित
है और अधिक बाधित है।

\textbf{4.} $\qty{108.0}{g/mol}$; $\qty{10.0}{g}/\qty{108.0}{g/mol} = \qty{0.0926}{mol}$।

\textbf{5.} हाइड्रैज़ोन $\ce{C9H12N2}$, \qty{148.0}{g/mol}: $\qty{0.0926}{mol} \times \qty{148.0}{g/mol} = \qty{13.7}{g}$।

\textbf{6.} $\ce{C6H5-NH-NH-C(CH3)=CH2}$।

\textbf{7.} \textit{ऑर्थो} वलय-कार्बन, \textit{इप्सो} कार्बन, दोनों नाइट्रोजन, हाइड्रैज़ोन कार्बन
और अंतस्थ $\ce{CH2}$।

\textbf{8.} N--N आबंध टूटता है; \textit{ऑर्थो} कार्बन और $\ce{CH2}$ के बीच एक C--C आबंध बनता है।

\textbf{9.} छह इलेक्ट्रॉन, सभी घटक \omterm{def:b3:pericyclic:faciality}{समफलकीय}: तापीय रूप से अनुमत।

\textbf{10.} लुइस अम्ल के रूप में यह चलावयवन को उत्प्रेरित करता है, एक नाइट्रोजन से बँधकर
N--N आबंध को दुर्बल करता है, और अमोनिया के निकलने में सहायता करता है।

\textbf{11.} \textit{ऑर्थो} कार्बन अब $sp^3$ हो गया है, एक हाइड्रोजन के साथ: उस प्रोटॉन का निकास
बेंज़ीन वलय को पुनः स्थापित करता है, जो एक प्रबल प्रेरक बल है।

\textbf{12.} बना हुआ ऐनिलीन $\ce{NH2}$ इमीन कार्बन पर आक्रमण करता है, जिससे
पंचसदस्यीय वलय बंद होता है (एक 2-ऐमीनोइंडोलीन)।

\textbf{13.} अमोनिया; इंडोल के ऐरोमैटिक पिरोल वलय का बनना इसे आगे बढ़ाता है।

\textbf{14.} $\ce{C9H12N2 -> C9H9N + NH3}$।

\textbf{15.} \qty{131.0}{g/mol}।

\textbf{16.} $\qty{0.0926}{mol} \times \qty{131.0}{g/mol} = \qty{12.1}{g}$।

\textbf{17.} C3 पर, जो इंडोल की सबसे नाभिकरागी स्थिति है और यहाँ मुक्त है; वहाँ आक्रमण
मध्यवर्ती में बेंज़ीन वलय को अक्षुण्ण रखता है।

\textbf{18.} $\ce{C6H5NHNH-C(=CH2)-CH2CH3}$ (अंतस्थ) और $\ce{C6H5NHNH-C(CH3)=CH-CH3}$ (आंतरिक)।

\textbf{19.} 2-एथिलइंडोल और 2,3-डाइमेथिलइंडोल।

\textbf{20.} आंतरिक वाला अधिक प्रतिस्थापित है; प्रबल अम्ल अंतस्थ वाले का
अंश बढ़ाते हैं।

\textbf{21.} 2,3-डाइमेथिलइंडोल, अधिक स्थायी, अधिक प्रतिस्थापित
ईन-हाइड्रैज़ीन के माध्यम से।

\textbf{22.} ऐसीटैल्डिहाइड स्वयं से संघनित हो जाता है और ईन-हाइड्रैज़ीन मार्ग अल्प
लब्धि देता है; इसके स्थान पर इंडोल पाइरुविक अम्ल के हाइड्रैज़ोन से बनाया जाता है, जिससे
इंडोल-2-कार्बोक्सिलिक अम्ल मिलता है, जिसका विकार्बोक्सिलन किया जाता है।

\textbf{23.} स्थान-समावयवियों के मिश्रण को पृथक करना पड़ता है, जिसकी कीमत लब्धि से चुकानी पड़ती है;
कोई सममित कीटोन या ऐसा मार्ग जो स्थान-रसायन को निश्चित कर दे, वरीय होता है।

\textbf{24.} 10.0 g फ़ेनिलहाइड्रैज़ीन अधिक से अधिक \qty{12.1}{g} 2-मेथिलइंडोल दे सकता है।
\end{solution}
